% Source: physical PDF pp. 336--339 (printed pp. 313--316), from the
% Section 30.3 heading through the final paragraph immediately before
% Problems.
% Inventory: Eqs. (30.3.1)--(30.3.4), no unnumbered displays, citations,
% figures, tables, footnotes, or typographic dividers.
% Corrected “incoming” to “outgoing” for the E^* external lines.
% Uncertainties: none.

\subsection{Calculations with Supergraphs}

We now consider how to use the results of the previous sections to calculate
the quantum effective action \(\Gamma[S,S^*]\) for a set of classical
potential superfields \(S_n(x,\theta,\bar\theta)\) and their adjoints.
Following the
prescription discussed in Section 16.1, this may be defined in terms of a sum
over all connected one-particle-irreducible supergraphs, consisting of
vertices to which are attached directed internal and external lines. For each
external line of type \(n\) ending or beginning in a vertex labelled
\(x,\theta,\bar\theta\), we include a c-number factor
\(S_n(x,\theta,\bar\theta)\) or
\(S_n^*(x,\theta,\bar\theta)\), respectively (but no propagator).
A vertex labelled \(x,\theta,\bar\theta\) with \(N\) incoming or \(N\)
outgoing lines labelled
\(n_1,n_2,\ldots,n_N\) yields a factor equal to \(i\) times the coefficient
of the \(S_1S_2\cdots S_N\) term in the superpotential \(\widetilde f(S)\) or
times the complex conjugate of this coefficient, respectively. An internal
line of any type coming out of a vertex labelled \(x,\theta,\bar\theta\)
and going into a vertex labelled \(x',\theta',\bar\theta'\) yields a
propagator factor, given by
Eqs.~\eqref{eq:30.2.10} and \eqref{eq:30.2.17} as
\begin{equation}
  -\frac{i}{4}\delta^2(\theta-\theta')
    \delta^2(\bar\theta-\bar\theta')\Delta_F(x-x')\,.
  \tag{30.3.1}\label{eq:30.3.1}
\end{equation}
In addition, as shown by Eq.~\eqref{eq:30.1.7}, a superderivative
\(\bar D^2\) acts on the propagators or external line \(S\)-factors of all
but one of the internal or external lines coming into any vertex, and a
superderivative \(D^2\) acts on the propagators or external-line
\(S^*\)-factors of all but one of the lines coming out of any vertex.
The product of these factors must be integrated over all \(x\)s,
\(\theta\)s, and \(\bar\theta\)s;
the quantum effective action is the sum of these integrals for all
one-particle-irreducible graphs.

By integrating by parts in superspace, the \(D^2\) and/or
\(\bar D^2\) operators accompanying any one propagator (say, one that
connects vertices labelled \(x,\theta,\bar\theta\) and
\(x',\theta',\bar\theta'\)) may be moved to
other propagators or external line factors. This leaves the factor contributed
by this internal line proportional to
\(\delta^2(\theta-\theta')\delta^2(\bar\theta-\bar\theta')\).
Integrating over \(\theta'\) and \(\bar\theta'\) then eliminates these delta
functions, and lets us replace \(\theta',\bar\theta'\) with
\(\theta,\bar\theta\) everywhere else. (In cases where several internal
lines connect the same pair of vertices, we need to use
\(\left[
  \delta^2(\theta-\theta')
  \delta^2(\bar\theta-\bar\theta')
\right]^2=0\).
Continuing in this way, we wind up with a
single full \(d^4\theta\) integral, with all \(D\)s and \(\bar D\)s acting
on the external line factors \(S_n\) and \(S_n^*\). That is, though
not usually local in spatial coordinates, \(\Gamma[S,S^*]\) is
\emph{local in the fermionic coordinates}.

The structure of this functional is governed by its invariance under the
`gauge' transformations \eqref{eq:30.2.4}. This tells us that every one of
the external line factors \(S_n\) or \(S_n^*\) must be acted on by
an operator \(\bar D^2\) or \(D^2\), respectively, with two possible
exceptions. The exceptions are that a term with only incoming or only outgoing
external lines, in which all but one of the external line factors \(S_n\)
or \(S_n^*\) are acted on by \(\bar D^2\) or \(D^2\), respectively,
and by no other superderivatives, is still invariant under the transformation
\eqref{eq:30.2.4}, even though it cannot be expressed as a full
\(d^4\theta\) integral of a functional only of the \(\Phi_n\) and/or
\(\Phi_n^*\) alone. This is because the change in such an amplitude under this
transformation could only come from the change in the external line factor
\(S_n\) or \(S_n^*\) that is not acted on by an operator \(\bar D^2\)
or \(D^2\), and this change is eliminated if we use integration by parts
to move one of the other \(\bar D^2\) or \(D^2\) operators to act
on it. As we saw in Section 30.1, such a term in
\(\Gamma[S,S^*]\) is the
\(\mathcal{F}\)-term of a functional of the \(\Phi_n\) or of the
\(\Phi_n^*\), and thus makes a correction to the superpotential or to its
complex conjugate. Aside from these exceptions, every term in
\(\Gamma[S,S^*]\) may be written as a full \(d^4\theta\) integral
of a functional only of the \(\Phi_n\) and/or \(\Phi_n^*\), and therefore
makes a correction to the \(D\)-term part of the effective action.

Also note that a term in the quantum effective action in which all but one of
the external line factors \(S_n\) or \(S_n^*\) is acted on by
\(\bar D^2\) or \(D^2\), respectively, and also by additional
\(\bar D^2\) or \(D^2\) operators (such as in combinations like
\(\bar D^2D^2\bar D^2S_n\) or
\(D^2\bar D^2D^2S_n^*\)) may also be written by
integration by parts as a term in which one of the extra \(\bar D^2\) or
\(D^2\) acts on the previously undifferentiated external line
\(S_n\) or \(S_n^*\) factor. Such a term may therefore be expressed
as a full \(d^4\theta\) integral of a functional only of the \(\Phi_n\)
and/or \(\Phi_n^*\), and therefore just makes another correction to the
\(D\)-term part of the action. The only terms in
\(\Gamma[S,S^*]\) that
cannot be written in this way are those that have \(E\) incoming external
lines and only \(E-1\) operators \(\bar D^2\) operating on the \(S_n\)
external line factors, or \(E^*\) outgoing external lines and only
\(E^*-1\) operators \(D^2\) operating on the \(S_n^*\) external
line factors.
We can therefore tell whether a supergraph can make a contribution to the
superpotential or its adjoint by simply counting the number of
\(\bar D^2\) or \(D^2\) operators contributed by the supergraph
to the corresponding term in \(\Gamma[S,S^*]\).

Let us count these superderivatives. Consider a connected graph with:
\(V_n\) vertices with \(n\) lines coming in; \(V_n^*\) vertices with \(n\)
lines going out; \(I\) internal lines; \(E\) external incoming lines; and
\(E^*\) external outgoing lines. These numbers are related by
\begin{equation}
  I+E=\sum_n nV_n\,,
  \qquad
  I+E^*=\sum_n nV_n^*\,.
  \tag{30.3.2}\label{eq:30.3.2}
\end{equation}
The total numbers of \(\bar D^2\) and \(D^2\) operators are then
\begin{equation}
  N_{\bar D}
  =\sum_nV_n(n-1)
  =I+E-V
  =L+V^*+E-1\,,
  \tag{30.3.3}\label{eq:30.3.3}
\end{equation}
and
\begin{equation}
  N_D
  =\sum_nV_n^*(n-1)
  =I+E^*-V^*
  =L+V+E^*-1\,,
  \tag{30.3.4}\label{eq:30.3.4}
\end{equation}
where \(V=\sum_nV_n\) is the total number of vertices with incoming lines;
\(V^*=\sum_nV_n^*\) is the total number of vertices with outgoing lines; and
\(L=I-V-V^*+1\) is the number of loops. We see that a graph with any loops
will have \(N_{\bar D}\geq E\) and \(N_D\geq E^*\), so that there are at
least enough \(\bar D^2\) operators to convert all \(S_n\) into
\(\Phi_n=\bar D^2S_n\) and its derivatives, and enough
\(D^2\) operators to convert all \(S_n^*\) into
\(\Phi_n^*=D^2S_n^*\) and its derivatives.
\emph{Any graph with loops thus yields only a contribution to the
full \(d^4\theta\) integral---in other words the \(D\)-term---of a
functional of the left-chiral superfields \(\Phi_n\) and their adjoints.}

The only way to obtain a contribution to an \(\mathcal{F}\)-term or its
adjoint is to have just \(N_{\bar D}=E-1\) operators \(\bar D^2\) or just
\(N_D=E^*-1\) operators \(D^2\), respectively. According to
Eqs.~\eqref{eq:30.3.3} and \eqref{eq:30.3.4}, such a graph would have \(L=0\)
and \(V^*=0\) or \(V=0\), respectively. In other words, since we are
considering only one-particle-irreducible graphs, we can only get an
\(\mathcal{F}\)-term from a graph with a single vertex with incoming lines,
and we can only get its adjoint from a graph with a single vertex with
outgoing lines. This contribution is nothing but the integrated
\(\mathcal{F}\)-term of the original superpotential, or its adjoint. Thus we
see again that \emph{there is no finite or infinite renormalization of
\(\mathcal{F}\)-terms to any order of perturbation theory}.
